\subsection{复合积与挠积的连接同态}

命题 7. --- a) 设

\[
0 \to \mathrm{P} ^ {\prime} \xrightarrow {f} \mathrm{P} \xrightarrow {g} \mathrm{P} ^ {\prime \prime} \to 0\tag{8}
\]

是右 A-模的一个正合序列， \(\theta\in\mathrm{Ext}_{\mathrm{A}}^{1}\left(\mathrm{P}^{\prime\prime},\mathrm{P}^{\prime}\right)\) 为相应的类，M 为左 A-模。连接同态

\[
\delta_ {n} (\mathcal {E}, \mathrm{M}): \operatorname{Tor} _ {n} ^ {\mathrm{A}} \left(\mathrm{P} ^ {\prime \prime}, \mathrm{M}\right)\rightarrow \operatorname{Tor} _ {n - 1} ^ {\mathrm{A}} \left(\mathrm{P} ^ {\prime}, \mathrm{M}\right) \quad \text { is the map } \quad\gamma \mapsto \theta \circ \gamma .
\]

\begin{enumerate}
\def\labelenumi{\alph{enumi})}
\setcounter{enumi}{1}
\tightlist
\item
  设
\end{enumerate}

\[
0 \rightarrow \mathrm{M} ^ {\prime} \rightarrow \mathrm{M} \rightarrow \mathrm{M} ^ {\prime \prime} \rightarrow 0\tag{$(\mathcal{E}_1)$}
\]

是左 A-模的一个正合序列， \(\theta_{1} \in \operatorname{Ext}_{\mathbf{A}}^{1}(\mathbf{M}^{\prime \prime}, \mathbf{M}^{\prime})\) 为相应的类，P 为右 A-模。连接同态

\[
\delta_ {n} (\mathrm{P}, \mathcal {E} _ {1}): \operatorname{Tor} _ {n} ^ {\mathrm{A}} (\mathrm{P}, \mathrm{M} ^ {\prime \prime}) \rightarrow \operatorname{Tor} _ {n - 1} ^ {\mathrm{A}} (\mathrm{P}, \mathrm{M} ^ {\prime}) \quad \text { is the map } \quad\gamma \mapsto \theta_ {1} \circ \gamma .
\]

设 \(\gamma\in\mathrm{Tor}_{n}^{\mathrm{A}}(\mathrm{P}^{\prime\prime},\mathrm{M})\) 为闭链 \(z^{\prime\prime}\in\dot{\mathrm{Z}}_{n}(\mathrm{L}(\mathrm{P}^{\prime\prime})\otimes_{\mathrm{A}}\mathrm{L}(\mathrm{M}))\) 的类，并设

\[
\begin{array}{c} 0 \longrightarrow \mathrm{P} ^ {\prime} \xrightarrow {f} \mathrm{P} \xrightarrow {g} \mathrm{P} ^ {\prime \prime} \longrightarrow 0 \\ \uparrow_ {u _ {1}} \quad \uparrow_ {u _ {0}} \quad \uparrow_ {1} \\ \mathrm{L} _ {1} (\mathrm{P} ^ {\prime \prime}) \xrightarrow {d _ {1}} \mathrm{L} _ {0} (\mathrm{P} ^ {\prime \prime}) \xrightarrow {p _ {0} ^ {\prime \prime}} \mathrm{P} ^ {\prime \prime} \longrightarrow 0 \end{array}
\]

为一个交换图。我们将用 \(p' : \mathrm{L}(\mathrm{P}') \to \mathrm{P}'\) 和 \(p'' : \mathrm{L}(\mathrm{P}''') \to \mathrm{P}''\) 表示复形的典范态射。根据定义， \(\delta(\gamma) \in \mathrm{Tor}_{n-1}^{\mathrm{A}}(\mathrm{P}', \mathrm{M})\) 如下获得：选取 \(x \in \mathrm{P} \otimes \mathrm{L}_n(\mathrm{M})\) 使得 \((g \otimes 1)(x) = (p'' \otimes 1)(z'')\) ，而 \(\delta(\gamma)\) 是使得下式成立的循环 \(z' \in Z_{n-1}(\mathrm{L}(\mathrm{P}') \otimes \mathrm{L}(\mathrm{M}))\) 的类

\[
(f \otimes 1) (p ^ {\prime} \otimes 1) (z ^ {\prime}) = (1 \otimes d _ {n}) (x).
\]

对于 \(0 \leqslant i \leqslant n\) ，设 \(z_{i}^{\prime \prime}\) 为 \(z^{\prime \prime}\) 在 \(\mathrm{L}_{i}(\mathrm{P}^{\prime \prime}) \otimes \mathrm{L}_{n - i}(\mathrm{M})\) 中的分量；我们有

\[
0 = \mathrm{D} z ^ {\prime \prime} = \sum_ {i} \left(d _ {i} \otimes 1 + (- 1) ^ {i} \otimes d _ {n - i}\right) (z _ {i} ^ {\prime \prime}),
\]

\(\mathrm{donc}(d_i \otimes 1)(z_i'') = (-1)^i \otimes d_{n-i+1}(z_{i-1}'')\) 并且特别地

\[
\left(d _ {1} \otimes 1\right) \left(z _ {1} ^ {\prime \prime}\right) = - 1 \otimes d _ {n} \left(z _ {0} ^ {\prime \prime}\right).
\]

于是我们选择 \(x = (u_0 \otimes 1)(z_0'')\) ：我们确实有

\[
(g \otimes 1) (x) = \left(p _ {0} ^ {\prime \prime} \otimes 1\right) z _ {0} ^ {\prime \prime} = \left(p ^ {\prime \prime} \otimes 1\right) \left(z ^ {\prime \prime}\right).
\]

由于

\[
\begin{array}{r l} (1 \otimes d _ {n}) (x) & = (u _ {0} \otimes 1) (1 \otimes d _ {n}) (z _ {0} ^ {\prime \prime}) = - (u _ {0} \otimes 1) (d _ {1} \otimes 1) (z _ {1} ^ {\prime \prime}) \\ & = - (f \otimes 1) (u _ {1} \otimes 1) (z _ {1} ^ {\prime \prime}), \end{array}
\]

由此可知 \(\delta(\gamma)\) 是循环 \(z' \in Z_{n-1}(\mathbf{L}(\mathbf{P}') \otimes_{\mathbf{A}} \mathbf{L}(\mathbf{M}))\) （满足 \((p' \otimes 1)(z') = -(u_1 \otimes 1)(z_1'')\) ）的类。但根据定义，类 \(\theta\) 在同构 \(\operatorname{Ext}_A^1(\mathbf{P}', \mathbf{P}') \to \mathrm{H}^1(\operatorname{Homgr}_A(\mathbf{L}(\mathbf{P}''), \mathbf{P'}))\) 之下对应于态射 \(f: \mathbf{L}(\mathbf{P}') (1) \to \mathbf{P}'\) （由 \(-u_1\) 定义）的类，且乘积 \(\theta \circ \gamma\) 是循环

\[
\overline {{{z}}} ^ {\prime} \in Z _ {n - 1} (L (P ^ {\prime}) \otimes_ {A} L (M)) \quad \text { tels   que } \quad (p \otimes 1) (\overline {{{z}}} ^ {\prime}) = f (z ^ {\prime \prime}) = - (u _ {1} \otimes 1) (z _ {1} ^ {\prime \prime}),
\]

这完成了 \(a\) 的证明。断言 \(b\) 由 \(a\) 经交换同构得出。

推论 1. --- 设 \(0 \to P' \to P \to P'' \to 0\) 为右 A-模的正合序列， \(0 \to M' \to M \to M'' \to 0\) 为左 A-模的正合序列。则由连接同态复合而成的同态

\[
\operatorname{Tor} _ {n} ^ {\mathbf {A}} \left(\mathrm{P} ^ {\prime \prime}, \mathrm{M} ^ {\prime \prime}\right)\rightarrow \operatorname{Tor} _ {n - 1} ^ {\mathbf {A}} \left(\mathrm{P} ^ {\prime \prime}, \mathrm{M} ^ {\prime}\right)\rightarrow \operatorname{Tor} _ {n - 2} ^ {\mathbf {A}} \left(\mathrm{P} ^ {\prime}, \mathrm{M} ^ {\prime}\right)
\]

与

\[
\operatorname{Tor} _ {n} ^ {\mathbf {A}} \left(\mathrm{P} ^ {\prime \prime}, \mathrm{M} ^ {\prime \prime}\right)\rightarrow \operatorname{Tor} _ {n - 1} ^ {\mathbf {A}} \left(\mathrm{P} ^ {\prime}, \mathrm{M} ^ {\prime \prime}\right)\rightarrow \operatorname{Tor} _ {n - 2} ^ {\mathbf {A}} \left(\mathrm{P} ^ {\prime}, \mathrm{M} ^ {\prime}\right)
\]

是相反的。

事实上，若 \(\theta\) 与 \(\theta_{1}\) 是与给定正合序列相关联的类，且若 \(\gamma\in\mathrm{Tor}_{n}^{\mathbf{A}}(\mathbf{P}^{\prime\prime},\mathbf{M}^{\prime\prime})\) ，则 \(\gamma\) 的像分别为 \(\theta\circ(\theta_{1}\circ\gamma)\) 与 \(\theta_{1}\circ(\theta\circ\gamma)\) ，因此由命题 6 它们是相反的。

我们采用 X，第 127 页的记号，并考虑左A-模的序列(S)以及与正合序列(9)相关联的连接同态。

\[
\operatorname{Tor} _ {m} ^ {\mathbf {A}} \left(\mathrm{P}, \mathbf {K} _ {i - 1}\right)\rightarrow \operatorname{Tor} _ {m - 1} ^ {\mathbf {A}} \left(\mathrm{P}, \mathbf {K} _ {i}\right);
\]

我们通过复合迭代连接同态得出

\[
\partial_ {m} (\mathrm{P}, \mathscr {S}): \operatorname{Tor} _ {m} ^ {\mathrm{A}} (\mathrm{P}, \mathrm{M}) \rightarrow \operatorname{Tor} _ {m - n} ^ {\mathrm{A}} (\mathrm{P}, \mathrm{N}).
\]

然后由命题 7 和 X，第 118 页，命题 3：

推论 2. --- 若 \(\theta \in \operatorname{Ext}_{\mathbf{A}}^{n}(\mathbf{M}, \mathbf{N})\) 是与正合序列 \((\mathcal{S})\) 相关联的类，则我们有 \(\partial_{m}(\mathrm{P}, \mathcal{S})(\alpha) = \theta \circ \alpha\) ，它对所有 \(\alpha \in \operatorname{Tor}_{m}^{\mathbf{A}}(\mathrm{P}, \mathbf{M})\) .

推论 3. --- 若所有模 \(\mathbb{R}_i\) , \(i = 1, \ldots, n\) 都是平坦的，则映射 \(\alpha \mapsto \theta \circ \alpha\) 从 \(\operatorname{Tor}_{m+n}^{\mathbf{A}}(\mathbf{P}, \mathbf{M})\) 到 \(\operatorname{Tor}_m^{\mathbf{A}}(\mathbf{P}, \mathbf{N})\) 对每个右 A-模 P 和每个整数 \(m > 0\) 都是双射的 .

这由推论 2 和正合序列得出。

\[
\operatorname{Tor} _ {m + n - i + 1} ^ {\mathbf {A}} \left(\mathrm{P}, \mathrm{R} _ {i}\right)\rightarrow \operatorname{Tor} _ {m + n - i + 1} ^ {\mathbf {A}} \left(\mathrm{P}, \mathrm{K} _ {i - 1}\right) \xrightarrow {\hat {c}} \operatorname{Tor} _ {m + n - i} ^ {\mathbf {A}} \left(\mathrm{P}, \mathrm{K} _ {i}\right)\rightarrow \operatorname{Tor} _ {m + n - i} ^ {\mathbf {A}} \left(\mathrm{P}, \mathrm{R} _ {i}\right)
\]

其中两端的项按假设为零。

同样地，若

\[
0 \rightarrow \mathrm{Q} \rightarrow \mathrm{S} _ {n} \rightarrow \mathrm{S} _ {n - 1} \rightarrow \dots \rightarrow \mathrm{S} _ {1} \rightarrow \mathrm{P} \rightarrow 0\tag{$(\mathcal{S}_1)$}
\]

是右 A-模的正合序列，且 M 是左 A-模，则可定义迭代连接同态

\[
\partial^ {m} (\mathcal {S} _ {1}, \mathrm{M}): \operatorname{Tor} _ {m} ^ {\mathrm{A}} (\mathrm{P}, \mathrm{M}) \rightarrow \operatorname{Tor} _ {m - n} ^ {\mathrm{A}} (\mathrm{Q}, \mathrm{M})
\]

并且我们有：

推论 4. --- 若 \(\theta_{1} \in \operatorname{Ext}_{\mathbf{A}}^{n}(\mathbf{P}, \mathbf{Q})\) 是与正合序列 \((\mathcal{S}_{1})\) 相关联的类，则我们有 \(\partial^{m}(\mathcal{S}_{1}, \mathbf{M})(\alpha) = \theta_{1} \circ \alpha\) ，它对所有 \(\alpha \in \operatorname{Tor}_{m}^{\mathbf{A}}(\mathbf{P}, \mathbf{M})\) 成立 .

